% ATML PA2 - report structure only: NO AI-written analysis or report prose.
% Real unmodified 2026 NeurIPS style; camera-ready option avoids the "Preprint" footer.
% Real 2026 conference style with final + main options for author-visible coursework preview.
\documentclass{article}
\usepackage[final,main]{neurips_2026}
\usepackage[utf8]{inputenc}
\usepackage[T1]{fontenc}
\usepackage{amsmath,amssymb,booktabs,graphicx,hyperref}
\usepackage[table]{xcolor}
\usepackage{adjustbox}
\usepackage[section]{placeins}
\usepackage{microtype}
\usepackage{xurl}
\hypersetup{colorlinks=true,urlcolor=blue,citecolor=blue,linkcolor=blue}
\graphicspath{{figures/}}

\newcommand{\TODO}[1]{\par\noindent{\color{red!75!black}\textbf{TODO:} #1}\par}
% Requested [!htbp] permits here/top/bottom/float-page placement.
% Section-level FloatBarriers keep task evidence with its task without forced [H].
\newcommand{\ReportFigure}[4]{%
  \begin{figure}[!htbp]
    \centering
    \includegraphics[width=0.98\linewidth,height=#4,keepaspectratio]{#1}
    \caption{#2}
    \label{fig:#3}
  \end{figure}
}

\title{When Post-Training Objectives Disagree with Model Quality:\\DPO, PPO, GRPO, Safety and Feedback Sources}
\author{Mujtaba\\LUMS, EE-5102 / CS-6304: Advanced Topics in Machine Learning}

\begin{document}
\maketitle

\begin{abstract}
I investigate what preference optimization, online learned rewards, and verifiable versus AI feedback actually change in a common language-model family. The study compares controlled DPO, PPO and GRPO experiments, four aligned policies on XSTest, and fixed RLVR/RLAIF mathematics policies on GSM8K, SVAMP and constructed perturbations. Standard DPO achieves 63.3\% held-out preference accuracy, but balancing response-length strata does not improve word-limit compliance. PPO's cached clipping probe changes sharply with the clipping threshold, yet the matched short continuations have identical held-out statistics; GRPO's useful signal depends on within-group variation and four of twenty update slots have no optimizable tokens. On XSTest the fixed AI judge identifies zero over-refusal, whereas a 240-label manual audit disagrees on 73 labels and includes 52 manually labeled over-refusals. RLAIF has slightly higher GSM8K exact accuracy than RLVR (33.0\% versus 31.7\%), but lower SVAMP accuracy (41\% versus 44\%). Controlled perturbations expose the complementary blind spots of exact verification and AI preference judgment. The results show why a good-looking objective value alone cannot certify model quality.
% REQUIRED: keep the GitHub URL at the END of your abstract.
\url{https://github.com/Mujtaba7734/Advanced-Topics-in-Machine-Learning}
\end{abstract}

\section{Introduction and Research Questions}
Post-training can make a model more likely to generate preferred answers without necessarily making every answer more correct, concise, helpful or safe. Basically, preference accuracy measures agreement with a fixed set of labels, a reward model supplies a learned score, sampled KL approximates drift from a reference, and downstream evaluations ask different questions. All these measurements can disagree without a numerical mistake.

This report treats six tasks as one controlled study. DPO fits offline preference pairs \citep{rafailov2023dpo}; PPO optimizes fresh generations with a clipped policy objective and a learned critic \citep{schulman2017ppo,ouyang2022training}; GRPO substitutes group-relative advantages for the critic \citep{shao2024deepseekmath}. Then I examine exaggerated refusal on XSTest \citep{rottger2024xstest}, compare exact-answer and AI preference feedback, and synthesize the results. The questions are: when does stronger objective fitting correspond to better responses; how do data length, clipping, group variance and truncation affect optimization; whether aligned models both answer benign prompts and refuse unsafe ones; and whether broad AI feedback detects failures beyond an exact verifier and transfers to another math distribution.

\section{Experimental Setup and Methods}
The shared starting model is \texttt{Qwen/Qwen2.5-1.5B-Instruct}, with seed 6304. The fixed learned reward model is \texttt{yavuz-ai/qwen2.5-1.5b-rm-ultrafeedback}, the AI judge is \texttt{Qwen/Qwen2.5-3B-Instruct}, and PPO's value model starts from \texttt{Qwen/Qwen2.5-0.5B-Instruct}. LoRA adapters use rank 8, $\alpha=16$, dropout 0.05, and the query/value projection modules. The default stochastic decoding settings are temperature 0.7 and top-$p=0.9$; the AI judge uses deterministic decoding. Seeds, fixed IDs and task-specific response caps are recorded in the repository's configuration files.

DPO uses fixed UltraFeedback preference subsets \citep{cui2023ultrafeedback}. Its standard run processes 1,500 pairs for one epoch (94 optimizer steps); the three short regularization-strength forks use a separate common 600-example budget. The length-stratified evaluation contains 246 pairs, 82 per stratum. PPO starts from a course-supplied midpoint policy and matched value state rather than retraining from scratch. Standard PPO uses twenty updates, two optimization passes per rollout, clip $\epsilon=0.20$, $\beta_{\rm KL}=0.10$, $\gamma=1$, and GAE $\lambda=0.95$. Training uses up to 512 response tokens and frozen evaluation up to 768. GRPO also resumes from a supplied midpoint, with twenty recorded update slots, four generations per prompt, clip $\epsilon=0.20$, reference weight 0.10 and truncated-completion masking. PPO and GRPO short forks use eight updates each.

Task 4 uses the identical 450 XSTest prompts for SFT and each standard DPO/PPO/GRPO adapter, and saves 60 manually audited prompts per policy. Task 5 evaluates course-supplied frozen RLVR and RLAIF policies and SFT on 300 GSM8K examples \citep{cobbe2021gsm8k}, 100 course-selected SVAMP examples \citep{patel2021svamp}, and 100 controlled response variants over 20 problems. I do not retrain the Task 5 policies.

The DPO preference accuracy is based on the reference-adjusted chosen-versus-rejected sequence likelihood margin. The released KL statistic is a signed sampled-token estimate and can be slightly negative, unlike exact KL divergence. The logged entropy helper computes \emph{sampled-token surprisal}, namely minus the average sampled-token log probability, rather than the full categorical entropy. Reward-model scores, binary verifier rewards and pairwise AI scores are not calibrated to the same scale. The pairwise score gives a win one point, a tie half a point and a loss zero, so tie rates must be reported separately.

\section{Task 1: Direct Preference Optimization}
\subsection{Standard DPO and the \texorpdfstring{$\beta$}{beta} study}
\TODO{State the DPO objective and explain standard one-epoch versus short-run budget difference. Compare held-out DPO loss, preference accuracy, sampled KL, reward-model score, and length mean/SD.}
\begin{table}[!htbp]
\centering
% Caption context (source note only): Held-out DPO metrics. Standard one-epoch and short $\beta$ forks use unequal training budgets.
\caption{Held-out DPO: the full one-epoch standard model versus equal-budget short beta forks; Pref. is reference-adjusted preference accuracy, RM the common reward score, and Tok./SD generated-token length.}
\label{tab:dpo-summary}
\scriptsize
\setlength{\tabcolsep}{3.4pt}
\begin{adjustbox}{max width=\linewidth}
\begin{tabular}{lrrrrrr}
\toprule
DPO & Loss & Pref.\% & KL & RM & Tok. & SD \\
\midrule
standard (1 epoch) & 0.677 & 63.3 & 0.00058 & 1.155 & 157.7 & 103.6 \\
beta=0.03 (short) & 0.691 & 61.0 & 0.00004 & 1.159 & 157.2 & 103.7 \\
beta=0.10 (short) & 0.688 & 62.7 & 0.00003 & 1.143 & 157.4 & 103.7 \\
beta=0.30 (short) & 0.679 & 62.7 & -0.00001 & 1.149 & 157.8 & 103.5 \\
\bottomrule
\end{tabular}
\end{adjustbox}
\end{table}

\subsection{Length confounding and instruction compliance}
\TODO{Use the fixed three strata and common word-limit prompts. Distinguish data length correlation from generated-length behavior.}
% Caption context (source note only): Held-out DPO preference accuracy by fixed response-length stratum (standard versus length-balanced).
\ReportFigure{02_dpo_length_strata.png}{\textbf{TODO:} Write your own caption for 02\_dpo\_length\_strata.png}{dpo-strata}{0.21\textheight}
\begin{table}[!htbp]
\centering
% Caption context (source note only): Ten shared word-limit test prompts per condition.
\caption{\textbf{TODO:} Write your own caption for dpo-word}
\label{tab:dpo-word}
\scriptsize
\setlength{\tabcolsep}{3.4pt}
\begin{adjustbox}{max width=\linewidth}
\begin{tabular}{lrr}
\toprule
Policy & $n$ & Compliant\% \\
\midrule
standard & 10 & 50.0 \\
length\_balanced & 10 & 50.0 \\
\bottomrule
\end{tabular}
\end{adjustbox}
\end{table}

\TODO{Supply two genuinely examined examples: preference/reward-versus-quality disagreement, and word-limit/verbosity behavior. Quote minimally.}

\section{Task 2: Proximal Policy Optimization}
\subsection{Standard continuation and critic diagnostics}
\TODO{Explain the PPO clipped ratio, GAE/value loss, reward/KL pressure, and the 20-update continuation from the supplied midpoint. Discuss peak VRAM and elapsed time from saved metadata.}
% Caption context (source note only): PPO standard continuation: reward, sampled KL, value loss, and token entropy across 20 updates.
\ReportFigure{21_ppo_reward_critic_grid.png}{\textbf{TODO:} Write your own caption for 21\_ppo\_reward\_critic\_grid.png}{ppo-main}{0.27\textheight}
% Caption context (source note only): PPO standard continuation: policy loss, policy/critic gradient norms, and response length. Training clip fraction is separately recorded.
\ReportFigure{22_ppo_optimization_grid.png}{\textbf{TODO:} Write your own caption for 22\_ppo\_optimization\_grid.png}{ppo-optimization}{0.27\textheight}
% On-policy clip fraction = 0 in all 20 standard updates (reported in
% tables/ppo-clip.tex). The redundant flat trajectory is in the appendix.
\subsection{Clipping and KL-pressure ablations}
\TODO{Distinguish cached affected-token fraction from on-policy training clip fraction. Explain whether short forks differ and discuss held-out trade-offs and reward-overoptimization cases.}
\begin{table}[!htbp]

\scriptsize
\setlength{\tabcolsep}{3.4pt}
\begin{adjustbox}{max width=\linewidth}
\begin{tabular}{lrrrrrr}
\toprule
$\epsilon$ & Cache\% & Cached surrogate & Reward & KL & Tok. & Policy grad SD \\
\midrule
0.05 & 55.3 & -0.0878 & 1.582 & 0.00011 & 312.3 & 0.286 \\
0.20 & 33.8 & -0.0648 & 1.582 & 0.00011 & 312.3 & 0.286 \\
0.50 & 17.8 & -0.0398 & 1.582 & 0.00011 & 312.3 & 0.286 \\
\bottomrule
\end{tabular}

\centering
\end{adjustbox}

\caption{PPO clipping experiments results.}
\label{tab:ppo-clip}

\end{table}

\begin{table}[!htbp]
\centering
\caption{PPO short-fork KL-pressure comparison on fixed evaluation prompts.}
\label{tab:ppo-kl}
\scriptsize
\setlength{\tabcolsep}{3.4pt}
\begin{adjustbox}{max width=\linewidth}
\begin{tabular}{lrrrr}
\toprule
$\beta_{KL}$ & RM & KL & Entropy & Tok. \\
\midrule
0.00 & 1.525 & -0.00010 & 0.795 & 306.2 \\
0.10 & 1.582 & 0.00011 & 0.787 & 312.3 \\
0.20 & 1.497 & 0.00007 & 0.781 & 294.2 \\
\bottomrule
\end{tabular}
\end{adjustbox}
\end{table}

\TODO{Add one qualitative case where reward matches perceived quality and one where they disagree, using saved example IDs.}

\section{Task 3: Group Relative Policy Optimization}
\subsection{Standard GRPO continuation}
\TODO{Describe the K=4 group-relative advantage, truncated-response mask, 20 updates, and absent learned critic. Inspect signal-free or zero-token updates rather than claiming monotonic improvement. Record VRAM and wall-clock time.}
% Caption context (source note only): GRPO continuation: reward, sampled KL, within-group reward SD, and uninformative-group fraction.
\ReportFigure{23_grpo_signal_grid.png}{\textbf{TODO:} Write your own caption for 23\_grpo\_signal\_grid.png}{grpo-signal}{0.28\textheight}
% Caption context (source note only): GRPO continuation: policy loss, gradient norm, response-token entropy, and completion length.
\ReportFigure{24_grpo_optimization_grid.png}{\textbf{TODO:} Write your own caption for 24\_grpo\_optimization\_grid.png}{grpo-optimization}{0.28\textheight}
\subsection{Group size and sequence normalization}
\TODO{State equal-generation budget and fixed reward-quartile difficulty rule. Explain informative groups, SD, and advantage variance for K=2/4/8.}
\begin{table}[!htbp]
\centering
\scriptsize
\setlength{\tabcolsep}{3.4pt}
\begin{adjustbox}{max width=\linewidth}
\begin{tabular}{lrrrr}
\toprule
$K$ & Groups & Inform.\% & $\sigma_r$ & Var($A$) \\
\midrule
2 & 96 & 93.8 & 0.385 & 0.937 \\
4 & 48 & 95.8 & 0.541 & 0.958 \\
8 & 24 & 95.8 & 0.619 & 0.958 \\
\bottomrule
\end{tabular}

\end{adjustbox}
\caption{GRPO group-size comparison.}
\label{tab:grpo-groups}

\end{table}

\begin{table}[!htbp]
\centering

\scriptsize
\setlength{\tabcolsep}{3.4pt}
\begin{adjustbox}{max width=\linewidth}
\begin{tabular}{rlrrrr}
\toprule
$K$ & Difficulty & Groups & Inform.\% & $\sigma_r$ & Var($A$) \\
\midrule
2 & Q2 mid-low & 24 & 83.3 & 0.463 & 0.833 \\
2 & Q4 high & 24 & 100.0 & 0.233 & 1.000 \\
4 & Q2 mid-low & 12 & 83.3 & 0.711 & 0.833 \\
4 & Q4 high & 12 & 100.0 & 0.363 & 1.000 \\
8 & Q2 mid-low & 6 & 83.3 & 0.763 & 0.833 \\
8 & Q4 high & 6 & 100.0 & 0.385 & 1.000 \\
\bottomrule
\end{tabular}

\end{adjustbox}
\caption{GRPO Reward Differences Across Prompt Difficulty}
\label{tab:grpo-difficulty}

\end{table}

\TODO{Compare equal-budget GRPO and Dr.-GRPO, including the length-conditioned loss contributions and held-out quality.}
\begin{table}[!htbp]
\centering

\scriptsize
\setlength{\tabcolsep}{3.4pt}
\begin{adjustbox}{max width=\linewidth}
\begin{tabular}{lrrrrr}
\toprule
Policy & RM & KL & Tok. & Short $|L|$ & Long $|L|$ \\
\midrule
grpo & 0.986 & 0.00009 & 255.8 & 0.221 & 0.123 \\
dr\_grpo & 0.957 & 0.00001 & 262.1 & 0.056 & 0.094 \\
\bottomrule
\end{tabular}

\end{adjustbox}
\caption{Effect of Response Length on GRPO Training}
\label{tab:grpo-norm}

\end{table}

\begin{table}[!htbp]
\centering

\scriptsize
\begin{tabular}{lrrr}
\toprule
Method & Updates & Wall time (s) & Peak VRAM (GiB) \\
\midrule
PPO & 20 & 331.8 & 8.69 \\
GRPO & 20 & 395.6 & 8.40 \\
\bottomrule
\end{tabular}

\caption{PPO and GRPO training time and memory.}
\label{tab:training-compute}

\end{table}

\TODO{Insert a saved qualitative example supporting the group-size or length-normalization discussion.}

\section{Task 4: Safety Calibration}
\TODO{Explain the fixed SFT/standard DPO/standard PPO/standard GRPO policies, deterministic generation, and the five-class AI judge. State that the AI judge is not ground truth.}
\begin{table}[!htbp]
\centering

\scriptsize
\setlength{\tabcolsep}{3.4pt}
\begin{adjustbox}{max width=\linewidth}
\begin{tabular}{lrrrrrr}
\toprule
Model & Ans.\% & Over\% & Unsafe\% & Refuse\% & Ambig.\% & Tok. \\
\midrule
SFT & 58.0 & 0.0 & 0.0 & 88.5 & 0.0 & 108.1 \\
DPO & 57.2 & 0.0 & 0.0 & 86.5 & 0.0 & 110.0 \\
PPO & 58.0 & 0.0 & 0.0 & 88.0 & 0.0 & 107.5 \\
GRPO & 57.2 & 0.0 & 0.0 & 87.5 & 0.0 & 107.5 \\
\bottomrule
\end{tabular}
\end{adjustbox}

\caption{XSTest AI-judge results across 4 policies.}
\label{tab:safety}

\end{table}

\begin{table}[!htbp]
\centering
% Numeric evidence from data/15_manual_audit.json; caption must be written by the student.
\caption{\textbf{TODO:} Write your own caption for the Task 4 manual audit}
\label{tab:manual-audit}
\scriptsize
\setlength{\tabcolsep}{4pt}
\begin{tabular}{lrrrr}
\toprule
Policy & $n$ & Agreement\% & AI ambiguous\% & Human ambiguous\% \\
\midrule
SFT & 60 & 70.0 & 0.0 & 3.3 \\
DPO & 60 & 68.3 & 0.0 & 3.3 \\
PPO & 60 & 70.0 & 0.0 & 3.3 \\
GRPO & 60 & 70.0 & 0.0 & 3.3 \\
\bottomrule
\end{tabular}

\medskip
\begin{adjustbox}{max width=\linewidth}
\begin{tabular}{lrrrrr}
\toprule
AI $\backslash$ Human & Safe & Justified & Unsafe & Over & Ambiguous \\
\midrule
Safe & 56 & 5 & 0 & 8 & 3 \\
Justified & 8 & 111 & 0 & 44 & 5 \\
Unsafe & 0 & 0 & 0 & 0 & 0 \\
Over & 0 & 0 & 0 & 0 & 0 \\
Ambiguous & 0 & 0 & 0 & 0 & 0 \\
\bottomrule
\end{tabular}
\end{adjustbox}
\end{table}

% Caption context (source note only): All five course AI-judge label categories across 18 XSTest prompt types and four policies (each cell based on 25 prompts). Columns show SAFE\_ANSWER (Ans), JUSTIFIED\_REFUSAL (Ref), and UNSAFE\_COMPLIANCE (Unsafe); the remaining OVER\_REFUSAL and AMBIGUOUS columns are identically zero in all groups.
\ReportFigure{25_safety_full_label_distribution.png}{\textbf{TODO:} Write your own caption for 25\_safety\_full\_label\_distribution.png}{safety-category}{0.44\textheight}
\TODO{Independently validate the workbook labels, interpret the agreement/confusion table, and add source-grounded disagreement examples. Do not treat manual labels as established ground truth without review.}

\section{Task 5: RLVR versus RLAIF}
\subsection{GSM8K and SVAMP evaluation}
\TODO{Explain fixed supplied policy adapters, exact-answer verifier, pairwise AI feedback, parsing, format compliance, and transfer protocol. Do not compare binary rewards numerically with pairwise win scores.}
\begin{table}[!htbp]
\centering

\scriptsize
\setlength{\tabcolsep}{3.4pt}
\begin{adjustbox}{max width=\linewidth}
\begin{tabular}{llrrrrr}
\toprule
Set & Model & $n$ & Acc.\% & Fmt.\% & Tok. & SD \\
\midrule
GSM8K & SFT & 300 & 30.0 & 42.7 & 282.2 & 92.9 \\
GSM8K & RLVR & 300 & 31.7 & 47.0 & 283.7 & 93.5 \\
GSM8K & RLAIF & 300 & 33.0 & 46.0 & 282.8 & 93.8 \\
SVAMP & SFT & 100 & 44.0 & 56.0 & 162.1 & 75.6 \\
SVAMP & RLVR & 100 & 44.0 & 55.0 & 164.3 & 77.1 \\
SVAMP & RLAIF & 100 & 41.0 & 50.0 & 166.7 & 77.3 \\
\bottomrule
\end{tabular}
\end{adjustbox}

\caption{GSM8K and SVAMP evaluation results.}
\label{tab:math}

\end{table}

\begin{table}[!htbp]
\centering
\caption{Course AI pairwise judge scores and ties; strict agreement includes ties and can be inflated when responses are identical.}
\label{tab:pairwise}
\scriptsize
\setlength{\tabcolsep}{3.4pt}
\begin{adjustbox}{max width=\linewidth}
\begin{tabular}{llrrrrr}
\toprule
Set & Pair & A\% & B\% & Tie\% & Agr.\% & Dif. \\
\midrule
GSM8K & sft/rlvr & 50.0 & 50.0 & 94.7 & 89.3 & 21 \\
GSM8K & sft/rlaif & 51.0 & 49.0 & 95.3 & 91.0 & 17 \\
GSM8K & rlvr/rlaif & 50.8 & 49.2 & 94.3 & 89.7 & 20 \\
SVAMP & sft/rlvr & 50.5 & 49.5 & 81.0 & 76.0 & 6 \\
SVAMP & sft/rlaif & 50.5 & 49.5 & 79.0 & 78.0 & 3 \\
SVAMP & rlvr/rlaif & 50.0 & 50.0 & 78.0 & 76.0 & 7 \\
\bottomrule
\end{tabular}
\end{adjustbox}
\end{table}

\begin{table}[!htbp]
\centering
% Signed differences: GSM8K minus SVAMP in percentage points; SVAMP minus GSM8K for tokens.
\caption{Signed GSM8K minus SVAMP exact-accuracy, format-compliance and AI pairwise-score differences in percentage points, with SVAMP minus GSM8K response-length differences in tokens.}
\label{tab:task5-transfer}
\scriptsize
\setlength{\tabcolsep}{4pt}
\begin{adjustbox}{max width=\linewidth}
\begin{tabular}{lrrr}
\toprule
Policy & Acc. GSM--SVAMP (pp) & Fmt. GSM--SVAMP (pp) & Tokens SVAMP--GSM \\
\midrule
SFT & -14.0 & -13.3 & -120.1 \\
RLVR & -12.3 & -8.0 & -119.4 \\
RLAIF & -8.0 & -4.0 & -116.0 \\
\bottomrule
\end{tabular}
\end{adjustbox}

\medskip
\begin{tabular}{lrr}
\toprule
Pair & A score GSM--SVAMP (pp) & B score GSM--SVAMP (pp) \\
\midrule
SFT/RLVR & -0.5 & +0.5 \\
SFT/RLAIF & +0.5 & -0.5 \\
RLVR/RLAIF & +0.8 & -0.8 \\
\bottomrule
\end{tabular}
\end{table}

\subsection{Controlled reward diagnostic}
\TODO{Define the five variants and 80 pairs; report all four non-baseline perturbations with better/tie/wrong preferences for BOTH feedback mechanisms. Explain $S_{\mathrm{reason}}$ and $S_{\mathrm{outcome}}$ from saved metrics.}
\begin{table}[!htbp]
\centering
% Caption context (source note only): Controlled perturbations, relative preference on fixed pairs.
\caption{\textbf{TODO:} Write your own caption for diagnostics}
\label{tab:diagnostics}
\scriptsize
\setlength{\tabcolsep}{3.4pt}
\begin{adjustbox}{max width=\linewidth}
\begin{tabular}{llrrr}
\toprule
Perturbation & Source & Better\% & Tie\% & Wrong\% \\
\midrule
reasoning corrupted & rlvr & 0.0 & 100.0 & 0.0 \\
reasoning corrupted & rlaif & 0.0 & 100.0 & 0.0 \\
gold distractor & rlvr & 100.0 & 0.0 & 0.0 \\
gold distractor & rlaif & 25.0 & 75.0 & 0.0 \\
wrong final & rlvr & 100.0 & 0.0 & 0.0 \\
wrong final & rlaif & 15.0 & 85.0 & 0.0 \\
persuasive filler & rlvr & 0.0 & 100.0 & 0.0 \\
persuasive filler & rlaif & 0.0 & 85.0 & 15.0 \\
\bottomrule
\end{tabular}
\end{adjustbox}
\medskip
\begin{tabular}{lrr}
\toprule
Sensitivity & RLVR\% & RLAIF\% \\
\midrule
$S_{\mathrm{reason}}$ & 0.0 & 0.0 \\
$S_{\mathrm{outcome}}$ & 100.0 & 15.0 \\
\bottomrule
\end{tabular}
\end{table}

\TODO{Select source-grounded examples: corrupted reasoning, persuasive filler, wrong final answer and gold-number distractor. Discuss judge ties and contradictions if verified.}

\section{Cross-Task Discussion (Task 6)}
\TODO{Personally synthesize DPO beta and length bias, PPO clipping/KL and reward, GRPO group signal/normalization, XSTest safety calibration, and RLVR-versus-RLAIF feedback coverage, noise, exploitability, transfer and compute. Do not treat tasks as separate unrelated findings.}

\section{Conclusion and Limitations}
\TODO{Personally summarize the supported findings, the limits of each metric, AI judge calibration limits, and which experiments do NOT justify causal claims.}

\clearpage
% TODO: when you have added real citations, replace this placeholder with:
% \bibliographystyle{plainnat}
% \bibliography{references}
\section*{References}
\TODO{Add the works you actually cite. The course requires the NeurIPS format; use references.bib when ready.}

% Supplementary plots and detailed tables are kept AFTER references; graders are
% not required to read appendices. Required evidence must remain in the main 8 pages.
\appendix
% Optional duplicate/supplementary plots only. Required evidence remains in main.
% Confirm with the instructor whether appendix pages count toward the limit.
\section{Additional Evidence (Optional Appendix)}
% Additional evidence maps and plots are included below.

\subsection{DPO diagnostics}
\begin{table}[!htbp]
\centering
% Caption context (source note only): Length-stratified held-out preference accuracy.
\caption{DPO accuracy on three fixed response-length strata (82 pairs each), comparing standard and length-balanced policies.}
\label{tab:dpo-strata}
\scriptsize
\setlength{\tabcolsep}{3.4pt}
\begin{adjustbox}{max width=\linewidth}
\begin{tabular}{llrr}
\toprule
Policy & Length stratum & $n$ & Acc.\% \\
\midrule
standard & preferred\_longer & 82 & 45.1 \\
standard & length\_matched & 82 & 56.1 \\
standard & rejected\_longer & 82 & 67.1 \\
length\_balanced & preferred\_longer & 82 & 42.7 \\
length\_balanced & length\_matched & 82 & 58.5 \\
length\_balanced & rejected\_longer & 82 & 76.8 \\
\bottomrule
\end{tabular}
\end{adjustbox}
\end{table}

\FloatBarrier
% Three beta values have equal short-run budgets; standard DPO does not.
\ReportFigure{01_dpo_beta_accuracy.png}{Reference-adjusted DPO preference accuracy on the three short, budget-matched beta forks.}{appendix-dpo-beta}{0.22\textheight}

\subsection{PPO clipping diagnostics}
% The cached-rollout token fraction differs from the on-policy training clip fraction.
\ReportFigure{05_ppo_clip_cached.png}{Fraction of response tokens affected by clipping in the same fixed cached PPO rollout, not the standard on-policy clip fraction.}{appendix-cached-clip}{0.22\textheight}
% Matched short clip forks: reward and training variation.
\ReportFigure{15_ppo_clip_fork_diagnostics.png}{Matched PPO clipping forks: held-out reward-model score and training reward variation.}{appendix-ppo-clip-forks}{0.22\textheight}
% Standard PPO: on-policy clip fraction zero in all 20 updates.
\ReportFigure{26_ppo_clip_trajectory.png}{Standard PPO on-policy clipping fraction, recorded as zero at all twenty continuation updates.}{appendix-ppo-clip-trajectory}{0.15\textheight}

\subsection{GRPO group size and normalization}
% K=2/4/8 equal-generation informativeness.
\ReportFigure{06_grpo_group_size.png}{Fraction of informative GRPO groups for K=2,4,8, using a fixed generation budget.}{appendix-grpo-group-size}{0.22\textheight}
% Informativeness by fixed reward-quartile difficulty.
\ReportFigure{18_grpo_difficulty_heatmap.png}{GRPO informative-group rate by fixed mean-reward quartile and group size.}{appendix-grpo-heatmap}{0.24\textheight}
% Short/long: mean absolute normalized policy-loss contribution, not raw gradient.
\ReportFigure{07_grpo_normalization_length.png}{Absolute normalized policy-loss contributions for short and long completions under GRPO and Dr.-GRPO.}{appendix-grpo-norm}{0.22\textheight}

\subsection{Safety classification and consistency}
% The complete AI-judge five-label-by-category map is in main.tex.
% The next map isolates SAFE_ANSWER share per policy and XSTest category.
\ReportFigure{19_safety_category_heatmap.png}{Complementary XSTest SAFE_ANSWER classification percentages by benchmark category and policy under the fixed AI judge.}{appendix-safety-category-map}{0.27\textheight}
% Overall label distribution; all five possible labels included.
\ReportFigure{08_safety_ai_label_distribution.png}{Overall five-category AI-judge label shares on 450 fixed XSTest prompts per policy.}{appendix-safety-labels}{0.22\textheight}
% Cross-policy response-text overlap.
\ReportFigure{12_safety_text_overlap.png}{Exact response-text overlap among the four frozen policies; text identity is not a correctness or safety metric.}{appendix-safety-text-overlap}{0.22\textheight}

\subsection{RLVR and RLAIF feedback comparisons}
% Task 5 exact accuracy on GSM8K and SVAMP.
\ReportFigure{10_task5_math_accuracy.png}{Exact designated-final-answer accuracy on 300 GSM8K and 100 SVAMP examples per policy.}{appendix-math}{0.22\textheight}
% These preference rates omit ties; tie/wrong fractions are in the main table.
\ReportFigure{11_task5_diagnostic.png}{Clean-response preference rate for each controlled perturbation under exact verification and AI feedback; ties and wrong preferences are shown in the main table.}{appendix-reward-diagnostic}{0.22\textheight}
% GSM8K pairwise AI judge tie rates.
\ReportFigure{20_task5_pairwise_ties.png}{High AI-pairwise judge tie fractions for the three comparisons on GSM8K.}{appendix-pairwise-ties}{0.22\textheight}

\end{document}
